\documentclass[11pt]{article}
\usepackage[margin=1in]{geometry}
\usepackage{amsmath,amssymb,amsthm,mathtools}
\usepackage[hidelinks]{hyperref}

\newtheorem{theorem}{Theorem}
\newtheorem{proposition}[theorem]{Proposition}
\newtheorem{corollary}[theorem]{Corollary}
\newtheorem{lemma}[theorem]{Lemma}
\theoremstyle{definition}
\newtheorem{definition}[theorem]{Definition}
\newtheorem{assumption}[theorem]{Assumption}
\theoremstyle{remark}
\newtheorem{remark}[theorem]{Remark}

\newcommand{\eps}{\varepsilon}
\newcommand{\E}{\mathbb{E}}
\newcommand{\F}{\mathcal{F}}
\newcommand{\Bern}{\mathrm{Bern}}
\newcommand{\TV}{\mathrm{TV}}

\title{Beyond planted strings: theory for the VOUCH extension\\[2pt]
\large Working note, not for circulation}
\author{Quang-Vinh Dang}
\date{}

\begin{document}
\maketitle

\noindent This note collects the new results the extension rests on. Notation follows
the VOUCH paper: a committed cohort of $m$ pairs, pair contents $C_i=(c_i^0,c_i^1)$,
inclusion coins $b_i$, all remaining randomness of training and unlearning $\xi$, a
score $s$ in the declared class $\F$, signs
$z_i = \mathbf 1[s(M,c_i^{b_i}) > s(M,c_i^{1-b_i})]$ (ties split by a committed coin),
and the realised cohort advantage $\Delta = \tfrac2m\sum_i z_i - 1$. VOUCH certifies
$\Delta$ exactly (its Theorem~2) and reaches the coin-averaged advantage through its
Proposition~3. The four results below extend that ladder in three directions.

\section{Setting and the canary generator}

\begin{assumption}[Generator]\label{ass:gen}
The pair contents $C_1,\dots,C_m$ are i.i.d.\ draws from a generator distribution
$P$ on ordered pairs that is \emph{label-symmetric}: $(c^0,c^1)\sim P$ implies
$(c^1,c^0)\sim P$. The coins $b_i$ are i.i.d.\ $\Bern(\frac12)$, independent of $C$ and
of $\xi$; $\xi$ is independent of $C$.
\end{assumption}

Both properties are under the auditor's control: they describe the auditor's own
sampler, not the data or the model.

\section{Semantic twins keep the exact null}

VOUCH's secrets are random strings, which have no paraphrase. A semantic twin pair
instead states two incompatible facts about one fictitious entity: $c^0$ says the
entity has attribute value $v$, $c^1$ says it has value $v'$, with $\{v,v'\}$ drawn
from a value set. The concern is that semantics might break the symmetry that makes
the VOUCH null exact. It does not, provided the generator is label-symmetric.

\begin{proposition}[Exact null for semantic twins]\label{prop:semantic}
Under Assumption~\ref{ass:gen}, if the deployed model is conditionally independent of
the coins given the contents (exact unlearning), then for every score $s$ and every
step $t$ of a reveal order chosen independently of the signs,
$Z_{\pi(t)}\mid\F_{t-1}\sim\Bern(\frac12)$ exactly. This holds in particular for
\emph{belief scores} that compare the two values under the same question,
$s(M,c^{k}) = \log P_M(v_k\mid q)$ averaged over a committed set of paraphrased
questions $q$.
\end{proposition}

\begin{proof}
The proof of VOUCH Theorem~1 uses only that, given the unordered pair opened at step
$t$, the two scores are fixed numbers and the coin deciding which is the in-twin's is
a fair coin independent of everything revealed earlier and of the model. Neither the
semantics of $c^0,c^1$ nor the form of $s$ enters, as long as $s(M,\cdot)$ is a
function of the model and the record alone. A belief score of the stated form is such
a function: both twins share the question set and differ only in the value
evaluated.
\end{proof}

\begin{remark}[Power, not validity, is what semantics costs]
A base model may prefer one value a priori (``born in Paris'' over ``born in
Porto''). That preference enters both arms symmetrically and leaves the null exact,
but it adds variance to $D_i$ and costs power. The base-calibrated score
$s(M,c)-s(M_0,c)$, with $M_0$ the pre-fine-tuning model, removes it, and remains in
$\F$ because $M_0$ is fixed before the coins are drawn.
\end{remark}

\section{From the cohort to the generator population}

VOUCH Proposition~3 averages over the coins with the contents held fixed. The
auditor also controls the contents, through Assumption~\ref{ass:gen}, so the same
argument extends to them.

\begin{assumption}[Bounded influence]\label{ass:influence}
For fixed $\xi$, changing a single coin $b_i$ changes at most $\kappa$ of the $m$
signs, and replacing the contents $C_i$ of a single pair changes at most $\ell$ of
them.
\end{assumption}

Under exact unlearning $\kappa=\ell=1$: the model carries no trace of the forget set,
so a pair's coin or contents can only affect that pair's own sign.

\begin{theorem}[Transfer to the generator population]\label{thm:population}
Let $\Delta_\xi = \E_{b,C}[\Delta\mid\xi]$ be the expected advantage of this pipeline
run over fresh cohorts from $P$. Under Assumptions~\ref{ass:gen}
and~\ref{ass:influence}, for any $\alpha'\in(0,1)$, with probability at least
$1-\alpha'$ over $(b,C)$,
\[
|\Delta - \Delta_\xi| \;\le\; u_m(\alpha') \;:=\;
\sqrt{\frac{2(\kappa^2+\ell^2)\log(2/\alpha')}{m}} .
\]
Consequently a cohort certificate at tolerance $\eps$ and level $\alpha$ implies
$|\Delta_\xi| < \eps + u_m(\alpha')$ at level $\alpha+\alpha'$.
\end{theorem}

\begin{proof}
Given $\xi$, $\Delta$ is a function of the $2m$ independent coordinates
$(b_1,\dots,b_m,C_1,\dots,C_m)$. By Assumption~\ref{ass:influence}, changing one
coin changes $\Delta$ by at most $2\kappa/m$ and changing one pair's contents by at
most $2\ell/m$. McDiarmid's inequality gives
$\Pr[|\Delta-\Delta_\xi|\ge u]\le 2\exp\bigl(-2u^2/\sum_k c_k^2\bigr)$ with
$\sum_k c_k^2 = 4(\kappa^2+\ell^2)/m$, that is
$2\exp\bigl(-u^2m/(2(\kappa^2+\ell^2))\bigr)$; setting it to $\alpha'$ gives $u_m$.
The consequence follows from a union bound with the certificate's own error.
\end{proof}

\begin{remark}[What is still not reached]
$\Delta_\xi$ averages over cohorts but conditions on the training randomness $\xi$:
it is the advantage of \emph{this} deployed pipeline run on the generator
population. Averaging over $\xi$ would need several independent training runs and is
not attempted. For an auditor of a deployed model, $\Delta_\xi$ is the relevant
target. At $m=512$, $\alpha'=0.05$ and $\kappa=\ell=1$ the inflation is $0.17$,
against $0.12$ for the coin-only transfer of VOUCH Proposition~3.
\end{remark}

\section{From canaries to organic records, and the cost of special-casing}

Let $P_o$ be a distribution of \emph{organic} record pairs: two records of the kind a
user actually deletes, with a hypothetical fair coin choosing which one was trained
on. For a pair content $c$, write
\[
h_\xi(c) \;=\; \E\bigl[\,2Z_i - 1 \;\big|\; C_i = c,\ \xi\,\bigr],
\]
the expected sign contribution of a pair with contents $c$ placed in the cohort, the
expectation being over the coins and over the other pairs' contents. Then
$\Delta_\xi = \E_{c\sim P}[h_\xi(c)]$, and we define the \emph{organic advantage}
$\Delta^o_\xi = \E_{c\sim P_o}[h_\xi(c)]$: the advantage the same pipeline run would
show on an organic pair planted in the same way.

\begin{definition}[Distinguishing advantage]
For a class $\mathcal H$ of functions from pair contents to $[0,1]$,
$\gamma_{\mathcal H}(P,P_o)=\sup_{\phi\in\mathcal H}\bigl|\E_P\phi-\E_{P_o}\phi\bigr|$.
With $\mathcal H$ all measurable functions this is the total variation distance.
\end{definition}

\begin{theorem}[Special-casing requires detection]\label{thm:detect}
Let $\phi_\xi = (1+h_\xi)/2$, which takes values in $[0,1]$. Then
\[
|\Delta^o_\xi - \Delta_\xi| \;=\; 2\,\bigl|\E_{P}\phi_\xi - \E_{P_o}\phi_\xi\bigr|
\;\le\; 2\,\gamma_{\mathcal H}(P,P_o)
\]
for any class $\mathcal H$ containing $\phi_\xi$. Equivalently: if the pipeline treats
canaries and organic records differently enough to open a gap $g$ between their
advantages, then $\phi_\xi$ is a test that tells canary pairs from organic pairs with
advantage at least $g/2$.
\end{theorem}

\begin{proof}
$\Delta_\xi-\Delta^o_\xi = \E_P h_\xi - \E_{P_o}h_\xi = 2(\E_P\phi_\xi-\E_{P_o}\phi_\xi)$
by linearity, and the right-hand side is at most $2\gamma_{\mathcal H}$ whenever
$\phi_\xi\in\mathcal H$.
\end{proof}

\begin{remark}[How to read Theorem~\ref{thm:detect}]
The theorem is a reduction, not a free guarantee. $\phi_\xi$ is computable by anyone
who can run the provider's pipeline on a pair and score both twins, so the relevant
class $\mathcal H$ is ``tests the provider can implement''. It turns the
honest-but-verifiable assumption of VOUCH into a quantitative one: a provider who
special-cases the cohort must, in effect, operate a canary detector, and the gap it
can open is at most twice that detector's advantage. Three consequences follow.
(i) If the canary list is committed but hidden until after deployment, the provider
can only detect canaries from their contents, which is exactly what $\gamma$
measures. (ii) An auditor cannot certify $\gamma$ from data: any detector we train
gives a \emph{lower} bound on $\gamma$, never an upper bound. We therefore report
the strongest detectors we can build as evidence, and state the organic claim
conditionally on $\gamma$. (iii) Making $P$ close to $P_o$ (natural-language
canaries in the user's own format) shrinks $\gamma$ for every detector at once.
\end{remark}

\begin{corollary}[The full ladder]\label{cor:ladder}
Under the assumptions of VOUCH Theorem~2, Theorem~\ref{thm:population} and
Theorem~\ref{thm:detect}, if a two-sided $(\eps,\alpha)$ cohort certificate is issued,
then with probability at least $1-\alpha-\alpha'$,
\[
|\Delta^o_\xi| \;<\; \eps \;+\; u_m(\alpha') \;+\; 2\gamma_{\mathcal H}(P,P_o).
\]
Each term has a separate cause: the tolerance the auditor chose, the size of the
cohort, and how far the canaries are from the data they stand in for.
\end{corollary}

\section{How fast the cohort certificate arrives}

VOUCH Theorem~5 bounds the certification time only through refutation, which fires
late. Here we bound it through wealth growth. The key observation is that under a
uniform reveal order the betting edge of the without-replacement process is
\emph{deterministic}.

\begin{lemma}[Deterministic edge]\label{lem:edge}
Fix the upper certificate arm with boundary $p_0=\frac12+\frac\eps2$. Let the cohort
carry $j$ ones, $\hat p=j/m<p_0$, and $\delta=p_0-\hat p$. Under a uniform reveal
order, the recentred boundary $p_0(t)$ of VOUCH Theorem~2 and the conditional mean
$r_t=\E[Z_{\pi(t)}\mid\F_{t-1}]$ satisfy, for every $t$ and every history,
\[
g_t \;:=\; p_0(t) - r_t \;=\; \frac{m\delta}{m-t+1}.
\]
\end{lemma}

\begin{proof}
With $S_{t-1}$ ones revealed, $p_0(t)=(mp_0-S_{t-1})/(m-t+1)$ and
$r_t=(j-S_{t-1})/(m-t+1)$; subtract.
\end{proof}

The edge grows as the cohort is exhausted, which is the finite-population advantage
over an i.i.d.\ stream with the same mean, where the edge stays $\delta$.

\begin{theorem}[Certification time through wealth growth]\label{thm:growth}
Run the upper arm with a uniform mixture of $K$ constant stakes
$\Lambda=\{\lambda_1,\dots,\lambda_K\}\subset(0,\frac12]$, refuting the null outright
if $p_0(t)>1$. Let
$A_t=\sum_{u\le t}g_u=m\delta\,(H_m-H_{m-t})$ and
$G_t=\sum_{u\le t}g_u^2=m^2\delta^2\sum_{k=m-t+1}^{m}k^{-2}$, with $H_n$ the
harmonic numbers, and for $\lambda\in\Lambda$
\[
\Psi_t(\lambda)=\lambda A_t-\frac{\lambda^2(t/4+G_t)}{2(1-\lambda)^2},\quad
V_t(\lambda)=\frac{\lambda^2(t/4+G_t)}{(1-\lambda)^2},\quad
R_\lambda=\log\frac{1+\lambda}{1-\lambda}.
\]
Then the stopping time $\tau$ of the arm (reaching $1/\alpha$ or refuting) satisfies,
for every $t<m$,
\[
\Pr[\tau>t]\;\le\;\min_{\lambda\in\Lambda}
\exp\!\Bigl(-\frac{x_t(\lambda)^2}{2\bigl(V_t(\lambda)+R_\lambda x_t(\lambda)/3\bigr)}\Bigr),
\qquad x_t(\lambda)=\bigl(\Psi_t(\lambda)-\log(K/\alpha)\bigr)_+ .
\]
\end{theorem}

\begin{proof}
Fix $\lambda\in\Lambda$ and write
$X_u=\log\bigl(1+\lambda(p_0(u)-Z_{\pi(u)})\bigr)$. Before refutation
$p_0(u)\in[0,1]$, so $y_u=\lambda(p_0(u)-Z_{\pi(u)})\in[-\lambda,\lambda]$. By
Taylor's theorem with remainder, $\log(1+y)\ge y-y^2/(2(1-\lambda)^2)$ for
$y\ge-\lambda$, hence
$\E[X_u\mid\F_{u-1}]\ge\lambda g_u-\lambda^2\E[(p_0(u)-Z)^2\mid\F_{u-1}]/(2(1-\lambda)^2)$,
and $\E[(p_0(u)-Z)^2\mid\F_{u-1}]=r_u(1-r_u)+g_u^2\le\frac14+g_u^2$. Summing, the
compensator of $\log W^\lambda_t$ is at least $\Psi_t(\lambda)$, a deterministic
quantity by Lemma~\ref{lem:edge}. Also $|\log(1+y)|\le|y|/(1-\lambda)$, so the
conditional variances sum to at most $V_t(\lambda)$, and the centred increments are
bounded below by $-R_\lambda$. After refutation, replace $X_u$ by its compensator
lower bound so that the centred sequence stays a martingale difference sequence;
this changes nothing on the event that refutation has not occurred by $t$. Freedman's
inequality then gives
$\Pr[\log W^\lambda_t<\Psi_t(\lambda)-x]\le\exp(-x^2/(2(V_t+R_\lambda x/3)))$.
Finally, the mixture wealth is at least $W^\lambda_t/K$, so
$\{\tau>t\}\subseteq\{\log W^\lambda_t<\log(K/\alpha)\}$ on the event of no
refutation by $t$; take $x=x_t(\lambda)$ and minimise over $\Lambda$.
\end{proof}

\begin{remark}[Reading the bound]
Since $1/(m-u+1)\ge1/m$, $A_t\ge\delta t$, with equality only for an infinite cohort;
the factor $\rho_t=m(H_m-H_{m-t})/t\approx1+t/(2m)$ is the finite-population
speed-up. Ignoring $G_t$ and the factor $(1-\lambda)^{-2}$ and optimising $\lambda$,
the drift is about $2\delta^2t\rho_t^2$, so the time scale is
$t^\ast\approx\log(K/\alpha)/(2\delta^2\rho_{t^\ast}^2)$, against
$\log(1/\alpha)/\mathrm{KL}\approx\log(1/\alpha)/(2\delta^2)$ for an i.i.d.\ stream.
This explains the VOUCH observation that the cohort column of its power table sits
below the super-population information limit: the cohort answers an easier question,
and the edge grows as the cohort runs out. The lower arm is the mirror statement with
$\delta'=\hat p-(1-p_0)$, and the two-sided certificate needs both arms (union bound).
\end{remark}

\begin{remark}[Numerical check]
\texttt{extension/experiments/check\_cohort\_growth.py} compares the bound with the
empirical tail of the same mixture bettor over 2{,}000 random orders for
$(m,\eps)\in\{(512,0.2),(1024,0.2),(1024,0.1),(2048,0.1)\}$ at $\hat p=\frac12$. The
bound is never violated at any of the 32 checked times per setting. It is
conservative: the time at which it falls below one half is $2.2$ to $3.8$ times the
empirical median. The same bettor against the fixed boundary on an i.i.d.\ stream
with the same mean is slower, for example a median of $527$ against $370$ pairs at
$m=1024$, $\eps=0.1$.
\end{remark}

\section{What remains open}

\begin{itemize}
\item Theorem~\ref{thm:growth} covers a constant-stake mixture, not the adaptive
experts (ONS, KT) of the VOUCH engine. Including the constant-stake grid among the
engine's experts inherits the bound at the cost of $\log$ of the expert count.
\item Theorem~\ref{thm:detect} is conditional on $\gamma$, which data can only bound
from below. Closing that gap needs either a computational assumption on the
provider or a cryptographic attestation; neither is attempted here.
\item The organic advantage $\Delta^o_\xi$ is a counterfactual quantity: organic
records were not included by a coin. It is the advantage the pipeline \emph{would}
show on organic data under randomised inclusion, which is the standard definition
of membership leakage, but not an observed rate.
\end{itemize}

\end{document}
